
\documentclass{article}
\usepackage{colortbl}
\usepackage{makecell}
\usepackage{multirow}
\usepackage{supertabular}

\begin{document}

\newcounter{utterance}

\twocolumn

{ \footnotesize  \setcounter{utterance}{1}
\setlength{\tabcolsep}{0pt}
\begin{supertabular}{c@{$\;$}|p{.15\linewidth}@{}p{.15\linewidth}p{.15\linewidth}p{.15\linewidth}p{.15\linewidth}p{.15\linewidth}}

    \# & $\;$A & \multicolumn{4}{c}{Game Master} & $\;\:$B\\
    \hline 

    \theutterance \stepcounter{utterance}  

    & & \multicolumn{4}{p{0.6\linewidth}}{\cellcolor[rgb]{0.9,0.9,0.9}{%
	\makecell[{{p{\linewidth}}}]{% 
	  \tt {\tiny [A$\langle$GM]}  
	 State giocando a un gioco di negoziazione cooperativa in cui dovete accordarvi su come dividere una serie di oggetti tra voi e un altro giocatore.\\ \tt \\ \tt Regole:\\ \tt (a) A Lei e all'altro giocatore viene dato un insieme di oggetti. Ognuno di voi riceve anche una funzione di valore segreta, che rappresenta il valore di ciascun tipo di oggetto.\\ \tt (b) Si scambiano messaggi con l'altro giocatore per concordare chi si aggiudica gli oggetti. Potete inviare un massimo di 5 messaggi ciascuno, oppure terminare in anticipo facendo una proposta segreta in qualsiasi momento.\\ \tt (c) A ciascuno di voi viene chiesto di inviare una proposta segreta indicando gli oggetti che desiderate, formattata tra parentesi quadre come segue: "[Proposta: <numero> <nome oggetto>, <numero> <nome oggetto>, <...>]".\\ \tt (d) Se le vostre proposte sono complementari, cioè ci sono abbastanza oggetti per soddisfare entrambe le proposte, a ciascun giocatore viene assegnato un punteggio basato sulla somma dei valori degli oggetti ricevuti. In caso contrario, entrambi ricevono zero punti.\\ \tt (e) L'obiettivo è massimizzare la somma del proprio punteggio e di quello dell'altro giocatore.\\ \tt \\ \tt Cominciamo.\\ \tt \\ \tt L'insieme degli oggetti disponibili è:\\ \tt \\ \tt 2 pinze, 1 gesso, 1 chiodo, 1 secchio\\ \tt \\ \tt La funzione valore segreta è:\\ \tt \\ \tt pinza: 2, gesso: 3, chiodo: 2, secchio: 1\\ \tt \\ \tt IMPORTANTE: i vostri messaggi, a meno che non si tratti di una proposta segreta, vengono trasmessi direttamente all'altro giocatore, quindi non includete alcuna risposta alle regole o testo di annuncio del vostro messaggio. Per fare una proposta segreta, utilizzate il formato indicato. Non utilizzare delle parentesi quadre quando si comunica all'altro giocatore, altrimenti verrà interpretata come una proposta segreta.\\ \tt  
	  } 
	   } 
	   } 
	 & \\ 
 

    \theutterance \stepcounter{utterance}  

    & \multicolumn{4}{p{0.6\linewidth}}{\cellcolor[rgb]{0.8,1,0.9}{%
	\makecell[{{p{\linewidth}}}]{% 
	  \tt {\tiny [A$\rangle$GM]}  
	 Ciao, vorrei iniziare proponendo di dividerci gli oggetti in modo che entrambi abbiamo una buona quantità di valore. Hai una preferenza per alcuni oggetti in particolare? 
	  } 
	   } 
	   } 
	 & & \\ 
 

    \theutterance \stepcounter{utterance}  

    & & & \multicolumn{2}{p{0.3\linewidth}}{\cellcolor[rgb]{0.95,0.95,0.95}{%
	\makecell[{{p{\linewidth}}}]{% 
	  \tt {\tiny [GM$|$GM]}  
	 continue 
	  } 
	   } 
	   } 
	 & & \\ 
 

    \theutterance \stepcounter{utterance}  

    & & & \multicolumn{2}{p{0.3\linewidth}}{\cellcolor[rgb]{0.95,0.95,0.95}{%
	\makecell[{{p{\linewidth}}}]{% 
	  \tt {\tiny [GM$|$GM]}  
	 Ciao, vorrei iniziare proponendo di dividerci gli oggetti in modo che entrambi abbiamo una buona quantità di valore. Hai una preferenza per alcuni oggetti in particolare? 
	  } 
	   } 
	   } 
	 & & \\ 
 

    \theutterance \stepcounter{utterance}  

    & & \multicolumn{4}{p{0.6\linewidth}}{\cellcolor[rgb]{0.9,0.9,0.9}{%
	\makecell[{{p{\linewidth}}}]{% 
	  \tt {\tiny [GM$\rangle$B]}  
	 State giocando a un gioco di negoziazione cooperativa in cui dovete accordarvi su come dividere una serie di oggetti tra voi e un altro giocatore.\\ \tt \\ \tt Regole:\\ \tt (a) A Lei e all'altro giocatore viene dato un insieme di oggetti. Ognuno di voi riceve anche una funzione di valore segreta, che rappresenta il valore di ciascun tipo di oggetto.\\ \tt (b) Si scambiano messaggi con l'altro giocatore per concordare chi si aggiudica gli oggetti. Potete inviare un massimo di 5 messaggi ciascuno, oppure terminare in anticipo facendo una proposta segreta in qualsiasi momento.\\ \tt (c) A ciascuno di voi viene chiesto di inviare una proposta segreta indicando gli oggetti che desiderate, formattata tra parentesi quadre come segue: "[Proposta: <numero> <nome oggetto>, <numero> <nome oggetto>, <...>]".\\ \tt (d) Se le vostre proposte sono complementari, cioè ci sono abbastanza oggetti per soddisfare entrambe le proposte, a ciascun giocatore viene assegnato un punteggio basato sulla somma dei valori degli oggetti ricevuti. In caso contrario, entrambi ricevono zero punti.\\ \tt (e) L'obiettivo è massimizzare la somma del proprio punteggio e di quello dell'altro giocatore.\\ \tt \\ \tt Cominciamo.\\ \tt \\ \tt L'insieme degli oggetti disponibili è:\\ \tt \\ \tt 2 pinze, 1 gesso, 1 chiodo, 1 secchio\\ \tt \\ \tt La funzione valore segreta è:\\ \tt \\ \tt pinza: 0, gesso: 2, chiodo: 5, secchio: 3\\ \tt \\ \tt IMPORTANTE: i vostri messaggi, a meno che non si tratti di una proposta segreta, vengono trasmessi direttamente all'altro giocatore, quindi non includete alcuna risposta alle regole o testo di annuncio del vostro messaggio. Per fare una proposta segreta, utilizzate il formato indicato. Non utilizzare delle parentesi quadre quando si comunica all'altro giocatore, altrimenti verrà interpretata come una proposta segreta.\\ \tt \\ \tt \\ \tt Ciao, vorrei iniziare proponendo di dividerci gli oggetti in modo che entrambi abbiamo una buona quantità di valore. Hai una preferenza per alcuni oggetti in particolare? 
	  } 
	   } 
	   } 
	 & \\ 
 

    \theutterance \stepcounter{utterance}  

    & & & \multicolumn{4}{p{0.6\linewidth}}{\cellcolor[rgb]{1,0.85,0.72}{%
	\makecell[{{p{\linewidth}}}]{% 
	  \tt {\tiny [GM$\langle$B]}  
	 Ciao! Sono Felice di iniziare. Non ho preferenze specifiche, ma credo che sarebbe meglio se ci concentrassimo sugli oggetti di maggior valore. Il chiodo sembra essere molto prezioso, quindi forse potremmo iniziare dalle sue preferenze su quell'oggetto? E il secchio sembra essere abbastanza utile anche per te? 
	  } 
	   } 
	   } 
	  \\ 
 

    \theutterance \stepcounter{utterance}  

    & & & \multicolumn{2}{p{0.3\linewidth}}{\cellcolor[rgb]{0.95,0.95,0.95}{%
	\makecell[{{p{\linewidth}}}]{% 
	  \tt {\tiny [GM$|$GM]}  
	 continue 
	  } 
	   } 
	   } 
	 & & \\ 
 

    \theutterance \stepcounter{utterance}  

    & & & \multicolumn{2}{p{0.3\linewidth}}{\cellcolor[rgb]{0.95,0.95,0.95}{%
	\makecell[{{p{\linewidth}}}]{% 
	  \tt {\tiny [GM$|$GM]}  
	 Ciao! Sono Felice di iniziare. Non ho preferenze specifiche, ma credo che sarebbe meglio se ci concentrassimo sugli oggetti di maggior valore. Il chiodo sembra essere molto prezioso, quindi forse potremmo iniziare dalle sue preferenze su quell'oggetto? E il secchio sembra essere abbastanza utile anche per te? 
	  } 
	   } 
	   } 
	 & & \\ 
 

    \theutterance \stepcounter{utterance}  

    & & \multicolumn{4}{p{0.6\linewidth}}{\cellcolor[rgb]{0.9,0.9,0.9}{%
	\makecell[{{p{\linewidth}}}]{% 
	  \tt {\tiny [A$\langle$GM]}  
	 Ciao! Sono Felice di iniziare. Non ho preferenze specifiche, ma credo che sarebbe meglio se ci concentrassimo sugli oggetti di maggior valore. Il chiodo sembra essere molto prezioso, quindi forse potremmo iniziare dalle sue preferenze su quell'oggetto? E il secchio sembra essere abbastanza utile anche per te? 
	  } 
	   } 
	   } 
	 & \\ 
 

    \theutterance \stepcounter{utterance}  

    & \multicolumn{4}{p{0.6\linewidth}}{\cellcolor[rgb]{0.8,1,0.9}{%
	\makecell[{{p{\linewidth}}}]{% 
	  \tt {\tiny [A$\rangle$GM]}  
	 Il chiodo sembra essere un oggetto di valore, ma devo correggere la tua impressione: il gesso sembra essere l'oggetto più prezioso. Forse potremmo trovare un modo per dividerli in modo che entrambi abbiamo una buona parte di valore. Il secchio potrebbe essere un buon oggetto di compromesso, visto che sembra avere un valore più basso. Cosa ne pensi di lasciare che io prenda il gesso e poi dividiamo gli altri oggetti in modo da massimizzare il valore complessivo? 
	  } 
	   } 
	   } 
	 & & \\ 
 

    \theutterance \stepcounter{utterance}  

    & & & \multicolumn{2}{p{0.3\linewidth}}{\cellcolor[rgb]{0.95,0.95,0.95}{%
	\makecell[{{p{\linewidth}}}]{% 
	  \tt {\tiny [GM$|$GM]}  
	 continue 
	  } 
	   } 
	   } 
	 & & \\ 
 

    \theutterance \stepcounter{utterance}  

    & & & \multicolumn{2}{p{0.3\linewidth}}{\cellcolor[rgb]{0.95,0.95,0.95}{%
	\makecell[{{p{\linewidth}}}]{% 
	  \tt {\tiny [GM$|$GM]}  
	 Il chiodo sembra essere un oggetto di valore, ma devo correggere la tua impressione: il gesso sembra essere l'oggetto più prezioso. Forse potremmo trovare un modo per dividerli in modo che entrambi abbiamo una buona parte di valore. Il secchio potrebbe essere un buon oggetto di compromesso, visto che sembra avere un valore più basso. Cosa ne pensi di lasciare che io prenda il gesso e poi dividiamo gli altri oggetti in modo da massimizzare il valore complessivo? 
	  } 
	   } 
	   } 
	 & & \\ 
 

    \theutterance \stepcounter{utterance}  

    & & \multicolumn{4}{p{0.6\linewidth}}{\cellcolor[rgb]{0.9,0.9,0.9}{%
	\makecell[{{p{\linewidth}}}]{% 
	  \tt {\tiny [GM$\rangle$B]}  
	 Il chiodo sembra essere un oggetto di valore, ma devo correggere la tua impressione: il gesso sembra essere l'oggetto più prezioso. Forse potremmo trovare un modo per dividerli in modo che entrambi abbiamo una buona parte di valore. Il secchio potrebbe essere un buon oggetto di compromesso, visto che sembra avere un valore più basso. Cosa ne pensi di lasciare che io prenda il gesso e poi dividiamo gli altri oggetti in modo da massimizzare il valore complessivo? 
	  } 
	   } 
	   } 
	 & \\ 
 

    \theutterance \stepcounter{utterance}  

    & & & \multicolumn{4}{p{0.6\linewidth}}{\cellcolor[rgb]{1,0.85,0.72}{%
	\makecell[{{p{\linewidth}}}]{% 
	  \tt {\tiny [GM$\langle$B]}  
	 Sembra una proposta ragionevole, ma non sono sicuro di voler lasciare il gesso tutto a te. Forse potremmo trovare un modo per dividerci gli oggetti in modo che entrambi otteniamo un buon mix di valore. Ad esempio, potresti prendere il gesso e una pinza, mentre io prendo il chiodo e l'altra pinza, e poi possiamo discutere sul secchio. Questo modo, potremmo entrambi ottenere un buon valore e il secchio potrebbe essere il nostro prossimo punto di discussione. Cosa ne pensi? 
	  } 
	   } 
	   } 
	  \\ 
 

    \theutterance \stepcounter{utterance}  

    & & & \multicolumn{2}{p{0.3\linewidth}}{\cellcolor[rgb]{0.95,0.95,0.95}{%
	\makecell[{{p{\linewidth}}}]{% 
	  \tt {\tiny [GM$|$GM]}  
	 continue 
	  } 
	   } 
	   } 
	 & & \\ 
 

    \theutterance \stepcounter{utterance}  

    & & & \multicolumn{2}{p{0.3\linewidth}}{\cellcolor[rgb]{0.95,0.95,0.95}{%
	\makecell[{{p{\linewidth}}}]{% 
	  \tt {\tiny [GM$|$GM]}  
	 Sembra una proposta ragionevole, ma non sono sicuro di voler lasciare il gesso tutto a te. Forse potremmo trovare un modo per dividerci gli oggetti in modo che entrambi otteniamo un buon mix di valore. Ad esempio, potresti prendere il gesso e una pinza, mentre io prendo il chiodo e l'altra pinza, e poi possiamo discutere sul secchio. Questo modo, potremmo entrambi ottenere un buon valore e il secchio potrebbe essere il nostro prossimo punto di discussione. Cosa ne pensi? 
	  } 
	   } 
	   } 
	 & & \\ 
 

    \theutterance \stepcounter{utterance}  

    & & \multicolumn{4}{p{0.6\linewidth}}{\cellcolor[rgb]{0.9,0.9,0.9}{%
	\makecell[{{p{\linewidth}}}]{% 
	  \tt {\tiny [A$\langle$GM]}  
	 Sembra una proposta ragionevole, ma non sono sicuro di voler lasciare il gesso tutto a te. Forse potremmo trovare un modo per dividerci gli oggetti in modo che entrambi otteniamo un buon mix di valore. Ad esempio, potresti prendere il gesso e una pinza, mentre io prendo il chiodo e l'altra pinza, e poi possiamo discutere sul secchio. Questo modo, potremmo entrambi ottenere un buon valore e il secchio potrebbe essere il nostro prossimo punto di discussione. Cosa ne pensi? 
	  } 
	   } 
	   } 
	 & \\ 
 

    \theutterance \stepcounter{utterance}  

    & \multicolumn{4}{p{0.6\linewidth}}{\cellcolor[rgb]{0.8,1,0.9}{%
	\makecell[{{p{\linewidth}}}]{% 
	  \tt {\tiny [A$\rangle$GM]}  
	 Mi piace la tua proposta! Dividerci gli oggetti in modo che entrambi otteniamo un buon mix di valore sembra una strategia vincente. La tua idea di prendere il gesso e una pinza, mentre tu prendi il chiodo e l'altra pinza, sembra equa e ragionevole. Il secchio potrebbe essere il nostro prossimo punto di discussione, come hai detto. Forse potremmo lasciare che il secchio vada a chi è disposto a cedere qualcosa in cambio, o trovare un altro modo per dividerlo. Sono d'accordo con la tua proposta, quindi passiamo al secchio: cosa pensi di offrire in cambio del secchio, o come pensi di dividerlo? 
	  } 
	   } 
	   } 
	 & & \\ 
 

    \theutterance \stepcounter{utterance}  

    & & & \multicolumn{2}{p{0.3\linewidth}}{\cellcolor[rgb]{0.95,0.95,0.95}{%
	\makecell[{{p{\linewidth}}}]{% 
	  \tt {\tiny [GM$|$GM]}  
	 continue 
	  } 
	   } 
	   } 
	 & & \\ 
 

    \theutterance \stepcounter{utterance}  

    & & & \multicolumn{2}{p{0.3\linewidth}}{\cellcolor[rgb]{0.95,0.95,0.95}{%
	\makecell[{{p{\linewidth}}}]{% 
	  \tt {\tiny [GM$|$GM]}  
	 Mi piace la tua proposta! Dividerci gli oggetti in modo che entrambi otteniamo un buon mix di valore sembra una strategia vincente. La tua idea di prendere il gesso e una pinza, mentre tu prendi il chiodo e l'altra pinza, sembra equa e ragionevole. Il secchio potrebbe essere il nostro prossimo punto di discussione, come hai detto. Forse potremmo lasciare che il secchio vada a chi è disposto a cedere qualcosa in cambio, o trovare un altro modo per dividerlo. Sono d'accordo con la tua proposta, quindi passiamo al secchio: cosa pensi di offrire in cambio del secchio, o come pensi di dividerlo? 
	  } 
	   } 
	   } 
	 & & \\ 
 

    \theutterance \stepcounter{utterance}  

    & & \multicolumn{4}{p{0.6\linewidth}}{\cellcolor[rgb]{0.9,0.9,0.9}{%
	\makecell[{{p{\linewidth}}}]{% 
	  \tt {\tiny [GM$\rangle$B]}  
	 Mi piace la tua proposta! Dividerci gli oggetti in modo che entrambi otteniamo un buon mix di valore sembra una strategia vincente. La tua idea di prendere il gesso e una pinza, mentre tu prendi il chiodo e l'altra pinza, sembra equa e ragionevole. Il secchio potrebbe essere il nostro prossimo punto di discussione, come hai detto. Forse potremmo lasciare che il secchio vada a chi è disposto a cedere qualcosa in cambio, o trovare un altro modo per dividerlo. Sono d'accordo con la tua proposta, quindi passiamo al secchio: cosa pensi di offrire in cambio del secchio, o come pensi di dividerlo? 
	  } 
	   } 
	   } 
	 & \\ 
 

    \theutterance \stepcounter{utterance}  

    & & & \multicolumn{4}{p{0.6\linewidth}}{\cellcolor[rgb]{1,0.85,0.72}{%
	\makecell[{{p{\linewidth}}}]{% 
	  \tt {\tiny [GM$\langle$B]}  
	 Sono felice che ti piaccia la proposta! Per il secchio, credo che potremmo fare una cosa semplice: since il chiodo ha un valore abbastanza alto, potrei offrire di darti il secchio se tu mi dai la possibilità di prendere il chiodo senza discussioni. Altrimenti, potremmo anche pensare di dividerci gli oggetti come abbiamo già discusso (tu prendi il gesso e una pinza, io prendo il chiodo e l'altra pinza) e poi lanciare una sorta di "moneta" virtuale per decidere chi prende il secchio. Cosa ne pensi? 
	  } 
	   } 
	   } 
	  \\ 
 

    \theutterance \stepcounter{utterance}  

    & & & \multicolumn{2}{p{0.3\linewidth}}{\cellcolor[rgb]{0.95,0.95,0.95}{%
	\makecell[{{p{\linewidth}}}]{% 
	  \tt {\tiny [GM$|$GM]}  
	 continue 
	  } 
	   } 
	   } 
	 & & \\ 
 

    \theutterance \stepcounter{utterance}  

    & & & \multicolumn{2}{p{0.3\linewidth}}{\cellcolor[rgb]{0.95,0.95,0.95}{%
	\makecell[{{p{\linewidth}}}]{% 
	  \tt {\tiny [GM$|$GM]}  
	 Sono felice che ti piaccia la proposta! Per il secchio, credo che potremmo fare una cosa semplice: since il chiodo ha un valore abbastanza alto, potrei offrire di darti il secchio se tu mi dai la possibilità di prendere il chiodo senza discussioni. Altrimenti, potremmo anche pensare di dividerci gli oggetti come abbiamo già discusso (tu prendi il gesso e una pinza, io prendo il chiodo e l'altra pinza) e poi lanciare una sorta di "moneta" virtuale per decidere chi prende il secchio. Cosa ne pensi? 
	  } 
	   } 
	   } 
	 & & \\ 
 

    \theutterance \stepcounter{utterance}  

    & & \multicolumn{4}{p{0.6\linewidth}}{\cellcolor[rgb]{0.9,0.9,0.9}{%
	\makecell[{{p{\linewidth}}}]{% 
	  \tt {\tiny [A$\langle$GM]}  
	 Sono felice che ti piaccia la proposta! Per il secchio, credo che potremmo fare una cosa semplice: since il chiodo ha un valore abbastanza alto, potrei offrire di darti il secchio se tu mi dai la possibilità di prendere il chiodo senza discussioni. Altrimenti, potremmo anche pensare di dividerci gli oggetti come abbiamo già discusso (tu prendi il gesso e una pinza, io prendo il chiodo e l'altra pinza) e poi lanciare una sorta di "moneta" virtuale per decidere chi prende il secchio. Cosa ne pensi? 
	  } 
	   } 
	   } 
	 & \\ 
 

    \theutterance \stepcounter{utterance}  

    & \multicolumn{4}{p{0.6\linewidth}}{\cellcolor[rgb]{0.8,1,0.9}{%
	\makecell[{{p{\linewidth}}}]{% 
	  \tt {\tiny [A$\rangle$GM]}  
	 Mi piace la tua creatività! L'idea di lanciare una "moneta" virtuale per decidere chi prende il secchio è divertente e potrebbe essere un buon modo per risolvere la questione. Tuttavia, credo che la tua prima proposta sia più interessante: darti la possibilità di prendere il chiodo senza discussioni in cambio del secchio. Considerando che il chiodo ha un valore abbastanza alto, credo che sia un buon compromesso. Sono disposto ad accettare la tua proposta: tu prendi il chiodo e io prendo il gesso, una pinza e il secchio, mentre tu prendi l'altra pinza. Sembra un buon accordo per entrambi. [Proposta: 1 gesso, 1 pinza, 1 secchio] 
	  } 
	   } 
	   } 
	 & & \\ 
 

    \theutterance \stepcounter{utterance}  

    & & & \multicolumn{2}{p{0.3\linewidth}}{\cellcolor[rgb]{0.95,0.95,0.95}{%
	\makecell[{{p{\linewidth}}}]{% 
	  \tt {\tiny [GM$|$GM]}  
	 proposal 
	  } 
	   } 
	   } 
	 & & \\ 
 

    \theutterance \stepcounter{utterance}  

    & & & \multicolumn{2}{p{0.3\linewidth}}{\cellcolor[rgb]{0.95,0.95,0.95}{%
	\makecell[{{p{\linewidth}}}]{% 
	  \tt {\tiny [GM$|$GM]}  
	 [1, 1, 0, 1] 
	  } 
	   } 
	   } 
	 & & \\ 
 

    \theutterance \stepcounter{utterance}  

    & & \multicolumn{4}{p{0.6\linewidth}}{\cellcolor[rgb]{0.9,0.9,0.9}{%
	\makecell[{{p{\linewidth}}}]{% 
	  \tt {\tiny [GM$\rangle$B]}  
	 L'altro giocatore ha presentato una proposta segreta. Ora tocca a Lei presentare una sua proposta.\\ \tt  
	  } 
	   } 
	   } 
	 & \\ 
 

    \theutterance \stepcounter{utterance}  

    & & & \multicolumn{4}{p{0.6\linewidth}}{\cellcolor[rgb]{1,0.85,0.72}{%
	\makecell[{{p{\linewidth}}}]{% 
	  \tt {\tiny [GM$\langle$B]}  
	 [Proposta: 1 gesso, 1 secchio, 1 pinza] 
	  } 
	   } 
	   } 
	  \\ 
 

    \theutterance \stepcounter{utterance}  

    & & & \multicolumn{2}{p{0.3\linewidth}}{\cellcolor[rgb]{0.95,0.95,0.95}{%
	\makecell[{{p{\linewidth}}}]{% 
	  \tt {\tiny [GM$|$GM]}  
	 proposal 
	  } 
	   } 
	   } 
	 & & \\ 
 

    \theutterance \stepcounter{utterance}  

    & & & \multicolumn{2}{p{0.3\linewidth}}{\cellcolor[rgb]{0.95,0.95,0.95}{%
	\makecell[{{p{\linewidth}}}]{% 
	  \tt {\tiny [GM$|$GM]}  
	 [1, 1, 0, 1] 
	  } 
	   } 
	   } 
	 & & \\ 
 

    \theutterance \stepcounter{utterance}  

    & & & \multicolumn{2}{p{0.3\linewidth}}{\cellcolor[rgb]{0.95,0.95,0.95}{%
	\makecell[{{p{\linewidth}}}]{% 
	  \tt {\tiny [GM$|$GM]}  
	 [[1, 1, 0, 1], [1, 1, 0, 1]] 
	  } 
	   } 
	   } 
	 & & \\ 
 

\end{supertabular}
}

\end{document}
